\subsection{Problem Statement}
\label{subsec:problemstellung}

Repeatedly initializing similar projects consumes effort outside product
development. Without a shared baseline, structure, tools, configuration, and
build processes differ; separate starting points increase maintenance effort
and may diverge technically.

The first baseline examined already had separate runtime components, tests,
builds, CI, and documented architectural boundaries. However, rules for size,
complexity, and selected dependencies were not operationalized through a
shared repository-wide gate. For example, the parent state
\texttt{0cbe1ba} of the later governance commit contained direct imports of
SQLAlchemy and \path{app.db} in a FastAPI router; \texttt{461fc751} moved the
readiness logic into a service and infrastructure wiring into the composition
root \parencite{kleiveist2026governance-commit}. The gap is therefore
observable not only in a target description but also in a corrected layer
violation.

In short change cycles, including AI-assisted development, a rule that exists
only in documentation can become visible late. This observation does not
support a general claim about the quality of AI-generated code; it establishes
a need for early, reproducible feedback in the development workflow examined.

The template must therefore combine standardization with flexibility: a common
basis should enable reuse without burdening projects with unnecessary
components. This case study examines how project profiles and optional
extensions address this trade-off and where adaptation is still required
\parencite{kleiveist2026architecture,kleiveist2026profiles}, and how the
governance added later respects this variability.
\beginnerinput{workspace/statement/100_einleitung/120}
